\chapter{Unit 6: Environmental Demography, Constitutional Law \& Movements}

\begin{formulabox}{Demographic Growth Models \& Doubling Time}
\begin{itemize}[leftmargin=1.2em]
    \item \textbf{Exponential Population Growth Model (Malthusian)}:
    \begin{align}
    \frac{dN}{dt} &= rN \\
    N(t) &= N_0 e^{rt}
    \end{align}
    where $N_0 = \text{initial population}$, $r = \text{intrinsic per-capita rate of increase } (b - d)$, $t = \text{time elapsed}$.
    \item \textbf{Rule of 70 (Doubling Time)}:
    \begin{equation}
    T_d = \frac{\ln(2)}{r} = \frac{0.69315}{r} \approx \frac{70}{\text{Growth Rate in Percentage } (r\%)}
    \end{equation}
    For example, if a population grows at $r = 2.0\%$ per annum, its doubling time is $T_d = 70 / 2 = 35\text{ years}$.
    \item \textbf{Logistic Population Growth Model (Verhulst-Pearl)}:
    \begin{equation}
    \frac{dN}{dt} = rN \left(1 - \frac{N}{K}\right)
    \end{equation}
    where $K = \text{environmental carrying capacity}$.
    \begin{itemize}
        \item When $N \ll K$, $dN/dt \approx rN$ (exponential phase).
        \item When $N = K/2$, growth rate $dN/dt$ reaches its absolute maximum (inflection point).
        \item When $N \to K$, $dN/dt \to 0$ (asymptotic plateau).
    \end{itemize}
    \item \textbf{Integrated Logistic Equation}:
    \begin{equation}
    N(t) = \frac{K}{1 + \left(\frac{K - N_0}{N_0}\right) e^{-rt}}
    \end{equation}
    \item \textbf{Demographic Transition Model (DTM)}:
    \begin{itemize}
        \item Stage 1 (Pre-Industrial): High birth rate, high death rate $\implies$ Slow growth.
        \item Stage 2 (Transitional): High birth rate, plunging death rate $\implies$ Rapid population explosion.
        \item Stage 3 (Industrial): Declining birth rate, low death rate $\implies$ Decelerating growth.
        \item Stage 4 (Post-Industrial): Low birth rate, low death rate $\implies$ Population stabilization or contraction.
    \end{itemize}
\end{itemize}
\end{formulabox}

\begin{lawbox}{Constitutional Environmental Provisions \& Environmental Statutes}
\begin{itemize}[leftmargin=1.2em]
    \item \textbf{Constitutional Triad}:
    \begin{itemize}
        \item \textbf{Article 21 (Fundamental Right)}: Expanded by the Supreme Court to guarantee the right to live in a healthy, pollution-free environment.
        \item \textbf{Article 48A (Directive Principle of State Policy - 42nd Amendment, 1976)}: Mandates the State to protect and improve the environment and safeguard forests and wildlife.
        \item \textbf{Article 51A(g) (Fundamental Duty - 42nd Amendment, 1976)}: Duty of every citizen of India to protect and improve the natural environment including forests, lakes, rivers, and wildlife.
    \end{itemize}
    \item \textbf{Chronology of Landmark Indian Environmental Acts}:
    \begin{center}
    \begin{tabularx}{\linewidth}{l c X}
    \toprule
    \textbf{Statute} & \textbf{Year} & \textbf{Core Regulatory Mandate} \\
    \midrule
    Wildlife (Protection) Act & 1972 & Hunting ban, National Parks, Sanctuaries, Schedule species. \\
    Water Act & 1974 & Created CPCB and SPCBs; consent to establish/operate outfalls. \\
    Forest (Conservation) Act & 1980 & Restricts diversion of forest land for non-forest purposes. \\
    Air Act & 1981 & Air pollution control areas; emission stack limits; noise included 1987. \\
    Environment (Protection) Act & 1986 & Umbrella legislation passed post-Bhopal; Section 3 power to Central Govt. \\
    Public Liability Insurance Act & 1991 & Mandatory no-fault liability insurance for hazardous industries. \\
    Biological Diversity Act & 2002 & Sovereign rights over genetic resources, NBA, Access \& Benefit Sharing. \\
    Forest Rights Act (FRA) & 2006 & Restitution of ancestral land rights to Scheduled Tribes \& forest dwellers. \\
    National Green Tribunal (NGT) Act & 2010 & Dedicated judicial fast-track environmental tribunal under Article 21. \\
    \bottomrule
    \end{tabularx}
    \end{center}
\end{lawbox}

\begin{lawbox}{Chronology of Major Environmental Movements in India}
\begin{itemize}[leftmargin=1.2em]
    \item \textbf{Bishnoi Movement (1730)}: Khejarli, Rajasthan. Amrita Devi and 363 Bishnois sacrificed lives protecting Khejri (\textit{Prosopis cineraria}) trees.
    \item \textbf{Chipko Movement (1973)}: Chamoli, Uttarakhand. Led by Gaura Devi, Sunderlal Bahuguna, Chandi Prasad Bhatt. Villagers hugged trees to prevent commercial clear-cutting. Resulted in 15-year Himalayan commercial logging ban (1980).
    \item \textbf{Silent Valley Movement (1978)}: Kerala. Led by Kerala Sastra Sahitya Parishad (KSSP) to save evergreen rainforests from a hydroelectric dam on the Kunthipuzha River. Led to creation of Silent Valley National Park (1984).
    \item \textbf{Appiko Movement (1983)}: Uttara Kannada, Karnataka. Led by Pandurang Hegde; southern Indian counterpart of Chipko.
    \item \textbf{Jungle Bachao Andolan (1982)}: Singhbhum, Bihar (now Jharkhand). Adivasi resistance against replacing native Sal forests with commercial Teak monocultures.
    \item \textbf{Narmada Bachao Andolan (1985--present)}: Led by Medha Patkar and Baba Amte opposing mega-dams (Sardar Sarovar) on the Narmada River, demanding rehabilitation and resettlement rights.
\end{itemize}
\end{lawbox}
